\GuideSpread{01. Put the zero at the center}{The conjecture is a statement about one real coordinate.}{01_centered.pdf}{
The nontrivial zeros lie in the strip \(0<\Re\rho<1\). Center the strip by writing \(z=\rho-1/2\). The Riemann hypothesis becomes \(\Re z=0\) for every zero, counted with multiplicity.

The picture locates the critical line and the two strip boundaries. Its dots are a geometric model, not measured zeta zeros. Moving the coordinate origin does not alter the conjecture. It makes a possible horizontal displacement visible as an exponential parameter.

Every universal criterion in this guide refers to the full zero multiset. A finite drawing can explain that criterion but cannot verify it.}{The unresolved coordinate is the horizontal displacement of every centered zero.}

\GuideSpread{02. Complete the zeta function}{The completion reveals reflection symmetry.}{02_completion.pdf}{
For \(\Re s>1\), the zeta function has both a Dirichlet series and an Euler product. Its completed form is
\[
\xi(s)=\tfrac12s(s-1)\pi^{-s/2}\Gamma(s/2)\zeta(s).
\]
It is entire and satisfies \(\xi(s)=\xi(1-s)\). On the critical line it is real.

The plot reuses 701 audited numerical values on \(0\le y\le35\). Multiplication by \(e^{\pi y/4}\) makes the rapidly decaying oscillations legible. This is a visual rescaling of the sampled function, not a new zero criterion. Nine higher precision controls and independent Hardy Z comparisons are documented in the source report.}{A sampled critical line curve shows the completion, not the location of all zeros.}

\GuideSpread{03. A zero carries its reflections}{Symmetry preserves any possible displacement.}{03_quartet.pdf}{
Conjugation and the functional equation send a centered zero to \(-z\), \(\overline z\), and \(-\overline z\). A hypothetical point \(\delta+i\gamma\) therefore travels with a quartet \(\pm\delta\pm i\gamma\), allowing coincidences when \(\delta=0\).

The displayed quartet is a model with \(\delta=0.2\) and \(\gamma=1.5\). These are chosen coordinates, not observed off line zeros. The arrows encode the exact symmetry of the actual zero multiset.

In a heat sum, each mode also acquires its complex heat phase. A quartet cannot generally be replaced by a phase free product of a cosine and a hyperbolic cosine.}{Reflection excludes a one sided escape from the critical line.}

\GuideSpread{04. Gaussian heat makes the sum converge}{Positive heat time damps large ordinates.}{04_damping.pdf}{
For \(t>0\), define
\[
\Theta_t(u)=\sum_\rho e^{t z_\rho^2+u z_\rho}.
\]
For real \(u\), every summand has magnitude at most \(e^{t/4+|u|/2}e^{-t\gamma^2}\). The unconditional zero count \(N(Y)=O(Y\log Y)\) makes the Gaussian tail summable, even with fixed polynomial weights.

The curves show the exact damping function for three chosen positive heat times. They do not approximate the complete zero sum. On compact translation sets and heat strips bounded away from zero, the same domination permits termwise differentiation.

Consequently the packet satisfies \(\partial_t\Theta=\partial_u^2\Theta\).}{The packet is defined without assuming the Riemann hypothesis.}

\GuideSpread{05. A small multiplier does not remove growth}{A displaced mode can be initially invisible.}{05_heat_growth.pdf}{
A hypothetical centered mode with real part \(\delta>0\) has amplitude proportional to \(e^{\delta u}\). Its Gaussian heat multiplier may be extremely small if its ordinate is large, but that multiplier is nonzero.

The figure is an envelope model. It omits the fixed heat multiplier and uses chosen displacements \(0\), \(0.1\), and \(0.2\). It displays amplitude, not squared energy, and it contains no sampled off line zero.

For the actual full packet, the monograph proves that the upper exponential growth rate of its running maximum is the supremum of the centered real parts. The proof uses Laplace poles, so it does not presume that different modes fail to cancel pointwise.}{Amplitude has exponent \(\delta\). Squared amplitude has exponent \(2\delta\).}

\GuideSpread{06. Boundedness meets a nonzero residue}{The Laplace transform turns displacement into a pole.}{06_laplace.pdf}{
If the full heat packet is bounded on \([0,\infty)\), its Laplace transform is holomorphic in \(\Re w>0\). Initially in \(\Re w>1/2\), absolute convergence identifies it with
\[
\sum_\rho\frac{e^{t z_\rho^2}}{w-z_\rho}.
\]
Each distinct centered zero contributes a nonzero residue, multiplied by its positive integer multiplicity. A pole with positive real part contradicts that holomorphic extension. Reflection then excludes negative real parts.

The pole positions in this diagram are hypothetical. Its shaded region is the exact domain supplied by boundedness. The manuscript proves the continuation and normal convergence needed by the argument.}{For any one positive heat time, boundedness on the full half line is equivalent to RH.}

\GuideSpread{07. Local energy retains the same obstruction}{A unit translation band still detects exponential escape.}{07_energy.pdf}{
Apply \(L=1/4-\partial_u^2\) to the packet, and call the result \(\widetilde\Theta\). Fix \(0<t_0<T\). On each unit band define
\[
E_J=\int_{t_0}^{T}\int_J^{J+1}|\widetilde\Theta_t(u)|^2\,du\,dt.
\]
The manuscript proves that RH is equivalent to \(\log(1+E_J)/J\to0\). Under RH, absolute Gaussian summation gives the stronger bound \(E_J=O(1)\).

The figure shows constant and exponential model energies. The red curve has rate \(2\delta\), as appropriate for a squared quantity. These chosen profiles illustrate the limit and are not computed energies of the actual zero packet.}{Subexponential growth controls every distant unit band, not a finite selection.}

\GuideSpread{08. The explicit formula has three channels}{Zeros equal a pole channel plus an analytic channel minus primes.}{08_channels.pdf}{
The Gaussian explicit formula gives the exact identity
\[
\Theta_t(u)=2e^{t/4}\cosh(u/2)+\mathfrak A_t(u)-\mathfrak P_t(u).
\]
The archimedean term is a Gaussian integral against \(-\log\pi+\Re\psi(1/4+iv/2)\). The prime term sums both translated Gaussians at \(\pm\log n\), weighted by \(\Lambda(n)/\sqrt n\), with prefactor \(1/(2\sqrt{\pi t})\).

The diagram displays the exact signs and normalization. Its blocks have no numerical heights. Large separate channels can cancel. An upper bound for the positive prime packet alone does not control their difference at the RH scale. The monograph states the Fourier convention and derives the constants.}{An exact change of representation carries the obstruction into the difference of channels.}

\GuideSpread{09. Remove the explicit poles}{The annihilator preserves every nontrivial zero mode.}{09_annihilator.pdf}{
The operator \(L=1/4-\partial_u^2\) kills both exponentials \(e^{u/2}\) and \(e^{-u/2}\). It therefore removes the explicit hyperbolic cosine channel.

On a zero mode it multiplies by
\[
\tfrac14-z_\rho^2=\rho(1-\rho),
\]
which is nonzero for every nontrivial zero. Thus \(\widetilde\Theta=\widetilde{\mathfrak A}-\widetilde{\mathfrak P}\) still contains the full zero obstruction.

On a positive heat strip, the digamma series supplies a uniform bound for the differentiated archimedean channel. Removing that analytic contribution does not supply the arithmetic cancellation in the remaining prime channel.}{The operator removes a known term and keeps the unknown one.}

\GuideSpread{10. Work in the physical Gaussian metric}{A coordinate repair starts with a norm convention.}{10_gaussian.pdf}{
The normalized Gaussian is
\[
g_{t,a}(u)=(2\pi t)^{-1/4}e^{-(u-a)^2/(4t)}.
\]
Its squared integral with respect to \(du\) is exactly one. The plot compares three heat widths at the same center. Wider Gaussians have smaller peaks because their norm is fixed.

Their overlap is \(\langle g_{t,a},g_{t,b}\rangle=e^{-(a-b)^2/(8t)}\). The log centers of neighboring large integers approach one another, so raw integer Gaussian coordinates become nearly dependent.

The numerical quadrature in the companion script checks the normalization and overlap on a large finite interval. The exact identities come from Gaussian integration.}{A stable coordinate system must preserve the physical norm.}

\GuideSpread{11. Orthogonal modes resolve nearby centers}{Hermite functions replace a nearly dependent frame.}{11_hermite.pdf}{
The scaled Hermite functions \(\psi_{r,t,c}\) form an orthonormal basis in the physical Gaussian metric. The first five modes are drawn at \(t=1/2\) and \(c=0\), where the scale becomes the standard Hermite scale.

The ground mode is positive. Higher modes contain more oscillation and encode displacement details. Their orthogonality is a property of the functions, not a statement that the arithmetic coefficients of a prime packet are small.

A finite trapezoid calculation checks the Gram matrix of the first five modes over a wide interval. It is a numerical normalization control. The infinite basis theorem is the classical Hermite result recorded in the source register.}{Orthogonality closes a conditioning problem. It preserves the arithmetic signal.}

\GuideSpread{12. A shifted Gaussian has Poisson mode mass}{The coherent coefficients are explicit.}{12_poisson.pdf}{
Write \(\sigma=(a-c)/(2\sqrt t)\) and \(\lambda=\sigma^2\). The exact expansion is
\[
g_{t,a}=e^{-\lambda/2}\sum_{r\ge0}\frac{\sigma^r}{\sqrt{r!}}\psi_{r,t,c}.
\]
The squared coefficients are \(e^{-\lambda}\lambda^r/r!\). The plot uses the declared example \(\lambda=6\). Retaining modes through \(R\) leaves exactly the Poisson mass beyond \(R\).

For this example, the discrete weights through mode 24 are computed by two independent recurrences and stored. The projection identity is exact. The bar heights are numerical evaluations of that identity. A displaced Gaussian needs more modes as its displacement increases.}{The truncation loss can be calculated before the arithmetic coefficients are known.}

\GuideSpread{13. Linear mode count can beat the aggregate loss}{A declared shell fixes the required truncation rate.}{13_truncation.pdf}{
For fixed \(C>0\) and \(t_0>0\), suppose active centers satisfy \(|a-c_J|\le C\sqrt J+O(1)\) and \(t\ge t_0\). Set \(\alpha=C^2/(4t_0)\). Then \(\lambda\le\alpha J+O(\sqrt J)\).

Choose \(q>\alpha\) and retain modes through \(R_J=\lceil qJ\rceil\). The norm tail has exponent \(-I(q\mid\alpha)J/2\), where \(I=q\log(q/\alpha)-q+\alpha\). The prime coefficient mass costs \(e^{J/2+o(J)}\), so \(I>1\) makes their combined compression loss decay exponentially.

The plotted rate function uses \(\alpha=1\). It is a parameter calculation, not an observed arithmetic cancellation. The strict inequalities and the positive heat width are essential.}{Only \(O(J)\) orthogonal coordinates are needed, but their values still need a signed estimate.}

\GuideSpread{14. The ground pairing is the field again}{Testing by a Gaussian advances the heat time.}{14_ground.pdf}{
Gaussian integration gives
\[
\int_{\mathbb R}g_{s,c}(u)\widetilde\Theta_t(u)\,du
=K_s\widetilde\Theta_{t+s}(c),
\]
where \(K_s=2\sqrt{\pi s}(2\pi s)^{-1/4}\).

The diagram traces the exact parameter changes: integrate against a normalized Gaussian, advance heat from \(t\) to \(t+s\), and evaluate at its center \(c\). It shows no measured packet values.

The pairing converges absolutely even without RH. It does not claim that the unweighted packet belongs to global \(L^2\). The lowest stable coordinate is another copy of the same field, so removing it without recovery would remove part of the problem.}{The coordinate repair retains a full strength ground channel.}

\GuideSpread{15. A classical sequence gives a second interface}{The binomial coefficient formula is finite and exact.}{15_ck.pdf}{
B\'aez Duarte's coefficient has two equal expressions:
\[
c_K=\sum_{j=0}^K\frac{(-1)^j\binom Kj}{\zeta(2j+2)}
=\sum_{n\ge1}\frac{\mu(n)}{n^2}(1-n^{-2})^K.
\]
The chart evaluates the finite formula for \(0\le K\le12\), using exact rational Bernoulli coefficients and an 80 digit value of \(\pi\). At \(K=0\), use the finite formula, or interpret the zeroth power as one. These thirteen numerical values do not test the asymptotic criterion.

RH is equivalent to \(|c_K|=O_\varepsilon(K^{-3/4+\varepsilon})\) for every positive \(\varepsilon\). Omitting that quantifier or replacing it by finitely many bounds changes the statement.}{The coefficient is classical. A cube decomposition will reorganize the same quantity.}

\GuideSpread{16. The four corners of the first prime cube}{For \(Q=6\), each corner has an exact divisor sign.}{16_cube6.pdf}{
Define \(w_K(n)=n^{-2}(1-n^{-2})^K\). For \(Q=2\cdot3\),
\[
\begin{aligned}
\Delta_6w_K(m)&=w_K(m)-w_K(2m)\\
&\quad-w_K(3m)+w_K(6m).
\end{aligned}
\]
Each coordinate decides whether to include one selected prime in the divisor \(b\). Corner signs are the Möbius signs of the divisors. The geometric square is an exact indexing diagram.

For the hand check \(K=1\),
\[
\Delta_6w_1(m)=\frac{2}{3m^2}-\frac{25}{27m^4}.
\]
The plotted bars use \(m=5\) and show the four rational summands. Their sum is \(17/675\), before multiplication by the remaining base sign \(\mu(5)=-1\).}{Alternating the corners does not remove the sign of the base.}

\GuideSpread{17. The eight corners of a larger cube}{For \(Q=30\), the third prime adds a new binary coordinate.}{17_cube30.pdf}{
The eight divisors of \(30=2\cdot3\cdot5\) index \(\Delta_{30}w_K(m)\). A corner containing an odd number of primes is negative. A corner containing an even number is positive.

The diagram uses two parallel squares for the third coordinate, making all eight labels readable. It is a combinatorial cube, not an embedded physical space.

At \(K=1\), exact rational summation gives
\[
\Delta_{30}w_1(m)=\frac{16}{25m^2}-\frac{208}{225m^4}.
\]
For \(m=7\), its value is \(6848/540225\). The full signed coefficient still attaches \(\mu(7)=-1\) to that cube.}{A larger cube reorganizes the weight and leaves an arithmetic sign on every rough base.}

\GuideSpread{18. The triangle majorant loses cancellation}{An upper bound can be too large while the signed sum is small.}{18_signed.pdf}{
For squarefree \(Q\), exact unique factorization gives
\[
c_K=\sum_{(m,Q)=1}\mu(m)\Delta_Qw_K(m).
\]
The full absolute majorant replaces each summand by \(|\mu(m)|\,|\Delta_Qw_K(m)|\). It is an upper bound for \(|c_K|\).

The plot uses a finite \(Q=6\), \(K=1\) example through \(m=100\). It compares the signed partial sum with the corresponding partial absolute mass. Both are finite quantities and both omit a tail. The source script records the exact rational totals.

The manuscript's asymptotic barrier lower bounds the full majorant. It supplies no lower bound for the signed amplitude.}{A lower bound on an upper bound diagnoses that upper bound, not the signed object.}

\GuideSpread{19. Squarefree support must survive the cube}{Zero Möbius coefficients contribute no mass.}{19_support.pdf}{
Only bases with \((m,Q)=1\) enter the rough sector. Among them, only squarefree bases have \(|\mu(m)|=1\). A base divisible by a prime square has \(\mu(m)=0\), regardless of the sign or size of its cube weight.

The table enumerates bases coprime to \(6\) through 49. The rows for 25 and 49 are explicitly zero. This is an exact finite support table, not a density approximation.

Counting all coprime integers would inflate the lower bound for the actual triangle majorant. The revised barrier proof instead counts squarefree coprime bases, using \(\mu(m)^2=\sum_{\ell^2\mid m}\mu(\ell)\).}{Coprime and squarefree are separate requirements. Both belong in the argument.}

\GuideSpread{20. Every cube has a positive far tail}{The divisor estimate is uniform in the selected primes.}{20_positive_tail.pdf}{
For \(K=M^2\) and \(m\ge RM\), the manuscript proves
\[
\Delta_Qw_{M^2}(m)\ge\frac{\kappa_R}{m^2},
\qquad\kappa_R=\zeta(2)^{-1}-\zeta(4)/R^2.
\]
Take \(R=3\), which makes the constant positive. The estimate follows from the Euler product and Bernoulli's inequality. It is uniform over squarefree \(Q\).

The numerical curves compare \(m^2\Delta_Qw_{900}(m)\) for \(Q=6\) and 30 over \(90\le m\le360\), against the same lower constant. These are declared finite checks of a proved inequality. They neither estimate the signed infinite sum nor verify RH.}{The weight becomes positive in the far tail. The remaining Möbius signs do not.}

\GuideSpread{21. Counting supplies enough supported tail bases}{The main term dominates under the stated cube growth condition.}{21_count.pdf}{
In \([RM,2RM]\), the count of squarefree bases coprime to \(Q\) is
\[
RM\frac{\varphi(Q)}Q\prod_{p\nmid Q}(1-p^{-2})
+O_R(2^d\sqrt M),
\]
where \(d\) is the number of primes in \(Q\). The product is at least \(1/\zeta(2)\), and \(\varphi(Q)/Q\ge1/(d+1)\).

The chart displays the scale model \(M/(d+1)\) versus \(2^d\sqrt M\), choosing \(d=\lfloor\sqrt{\log M}\rfloor\). The curves omit constants and are not exact base counts. This declared choice satisfies \(d=o(\log M)\), under which the error is eventually smaller than the main term.}{A density estimate becomes useful only after its support and error are controlled.}

\GuideSpread{22. The sign free bound misses the target rate}{The obstruction concerns one absolute value route.}{22_barrier.pdf}{
Combining supported base counts with the positive tail gives
\[
\mathcal A_Q(M)\ge\frac{\kappa_R}{8R\zeta(2)}
\frac{\varphi(Q)}{QM}
\ge\frac{\kappa_R}{8R\zeta(2)M(d+1)}.
\]
For \(d=o(\log M)\), this has scale \(M^{-1-o(1)}\). The RH criterion needs signed decay \(M^{-3/2+\varepsilon}\) for every positive \(\varepsilon\).

The rate curves are schematic powers with constants suppressed. The gap does not prove that \(|c_{M^2}|\) is large. It proves that the absolute majorant cannot provide the required family of upper bounds. The earlier interval and support errors are repaired in the companion proof.}{The missing improvement is cancellation, not another sign free estimate of the same majorant.}

\GuideSpread{23. Two rigorous interfaces, one open problem}{The dependency map names every proved and unproved step.}{23_dependencies.pdf}{
The heat and sequence interfaces each carry a full RH criterion. The heat route asks for subexponential local energy. The arithmetic route asks for the signed bound on \(c_K\). Each is equivalent to RH by the stated proofs and classical theorem.

Gaussian and Hermite identities provide exact representation and stable compression. Prime cubes provide exact grouping. The barrier shows why their absolute majorant is insufficient.

The diagram deliberately contains no direct new transfer arrow between the Möbius sum and the prime heat packet. The two interfaces meet through RH itself. The final signed rough base estimate is still open, even though the cube identity is exact.}{A completed representation is a rigorous interface to the gap, not a closure of it.}

\GuideSpread{24. Read the evidence at its actual strength}{A finite check and a universal implication have different jobs.}{24_checks.pdf}{
The guide's exact computations cover divisor signs, cube coefficients, Möbius support, finite grouping, and finite recurrence identities. Its numerical controls cover cached critical line values, Gaussian normalization, Hermite orthogonality, Poisson coefficients, and the stated tail examples.

Its chosen off line modes, energy profiles, counting scales, and barrier rates are labeled models. None are observations of an off line zero or measured asymptotic cancellation.

The expanded monograph contains the proofs and source boundary. The source bundle records figure inputs, software, bounds, hashes, failures, and reproduction commands. Historical Lean certificates remain historical. This guide makes no fresh Lean claim.

The question for a future estimate is concrete: does it preserve the actual signs and control the full observable at the required scale?}{RH remains open in this work. The identities, corrected barrier, and exact remaining estimate are explicit.}
